% New English translation from checked cote 25, site 101 (PDF page 102).
% AI-translated by GPT-6 Astra (OpenAI, ChatGPT Pro mode).
\scanpage{25/101}
\mspage{4}

% NOTE: the reviewed source records long diagonal pen strokes without treating them as deletions of the passages.
\var{\margin{\ensuremath{/}}}{\placedmargin{\ensuremath{/}}}{marginal note placed}{ega25_25_101-v0001}is isomorphic to an $\underline{O}_P(n)$. \ -- \ From the point of view of the write-up, it was
awkward to begin by considering an $\underline{L}$ of finite type. Begin directly
by considering \ $L = \Gamma(E^{r+1}-(0),\ q_*(\underline{L}))$ \ defining the Module
$\overset{\hand{\sim}}{L} \ = \unsure{q|i|4}_*(q^*(\underline{L}))$ \ , we know (Chap II) that \ $L$ is a graded Module that
determines precisely $\underline{L}$, then we show that $L$ is free\var{\inserted[typed]{of rank 1}}{ of rank 1}{insertion}{ega25_25_101-v0002} as a graded
Module \ by the reasoning indicated\var{\margin{\ensuremath{\swarrow}}}{\placedmargin{\ensuremath{\swarrow}}}{marginal note placed}{ega25_25_101-v0003}\unsure{\hand{:}|\var{\over{.}{:}}{:}{writing over}{ega25_25_101-v0004}}

b) \ General case. \ It follows from case a) \ by \ III 4.6.5., using
the relation \ $H^1(P^r_k,\underline{O}_{P^r_k}) = 0$ \ established in III 2. \ \ QED\var{\margin{\ensuremath{/}}}{\placedmargin{\ensuremath{/}}}{marginal note placed}{ega25_25_101-v0005}

Variant of \ 1.4. \ Let \ $\underline{Z}\,(S)$ \ be the set of locally constant
integer-valued functions on S; we define an evident homomorphism
\[ Z(S) \hand{\longrightarrow} \mathrm{Pic}(P) \ , \]
\var{\inserted[typed]{the}}{the}{insertion}{ega25_25_101-v0006} \var{\xout{\illegible{1}}}{}{typewriter strike-out}{ega25_25_101-v0007} $\underline{n} \hand{\in} Z(S)$ \ corresponding in eff\var{\xout{\unsure{ect|ect}}}{}{typewriter strike-out}{ega25_25_101-v0008} \var{\inserted[typed]{to the}}{to the}{insertion}{ega25_25_101-v0009} \var{\xout{a}}{}{typewriter strike-out}{ega25_25_101-v0010} partitions $(S_n)_{n\hand{\in}\underline{Z}}$ \ into disjoint open\unsure{s}
(some of which may be empty), \var{\xout{in such a way that}}{}{typewriter strike-out}{ega25_25_101-v0011} and \var{\xout{in}}{}{typewriter strike-out}{ega25_25_101-v0012} with such
a partition we associate the invertible Module $\underline{O}_P(\underline{n})$ whose \var{\inserted[typed]{restr to}}{restr to}{insertion}{ega25_25_101-v0013} \var{\xout{value on}}{}{typewriter strike-out}{ega25_25_101-v0014}
$f^{-1}(S_n)$ \ is $\underline{O}_P(n)$. \ We therefore find a variant
(x bis) \hfill $\mathrm{Pic}(S)\mathrm{x}\underline{Z}(S) \hand{\longrightarrow} \mathrm{Pic}(P)$ \ , \hfill
and \ a more general and more satisfactory statement than 1.4. asserts that
(under the conditions of 1.1.) this is a bijection. (NB When S $\hand{\neq \emptyset}$
\var{\inserted[typed]{resp. connected}}{resp. connected}{insertion}{ega25_25_101-v0015}
then the \ canonical map $\underline{Z} \hand{\to} \underline{Z}(S)$ \ associating with every $n \hand{\in} \underline{Z}$ \ the constant
function of value n is injective resp. surjective, \ and we recover .
formally 1.4., which should therefore be stated in the more
general form I have just made explicit. \ Also point out in a remark that
in the language of the Picard scheme to be introduced in Chap V, 1.1.
(in the preceding equivalent form) is stated simply \ by saying that
the canonical homomorphism
\[ \underline{Z}_S \hand{\longrightarrow} \underline{\mathrm{Pic}}_{P/S} \]
from the constant group scheme \var{\xout{\unsure{on} \underline{Z} S defined by}}{}{typewriter strike-out}{ega25_25_101-v0016} $\underline{Z}$ over S into the scheme\var{\margin{\ensuremath{\backslash}}}{\placedmargin{\ensuremath{\backslash}}}{marginal note placed}{ega25_25_101-v0017}
\tnote{25-101-1}{The source's q-subscript-star in the first section formula, and q, i, or 4 in the next formula, are preserved. The French canceled endings et or etu in en effet give the same English ending ect in in effect; both uncertain source candidates are retained in this note, and the corrected text keeps the visible fragment in eff. Bis means the second version of the formula label; the label is retained. The isolated period after recover and the open NB parenthesis also follow the source.}

\var{\margin{101}}{\placedmargin{101}}{marginal note placed}{ega25_25_101-v0018}
